\documentclass[8pt]{beamer}
%aspectratio=1610

\usepackage{amssymb}
\usepackage{amsfonts}
\usepackage{amsmath}
\usepackage{amsthm}

\usepackage{mathtools}
\usepackage[T1]{fontenc}
\usepackage{booktabs}
\usepackage{tikz-cd}
\usepackage{stmaryrd}

\usefonttheme{serif}
\usepackage{newpxtext,newpxmath}

\usetheme{Berlin}
\setbeamertemplate{theorems}[numbered]
\setbeamertemplate{navigation symbols}{}

\definecolor{MySoftPink}{HTML}{F8BBD0}
\definecolor{MyDarkTeal}{HTML}{00838F} 
\definecolor{MyVividOrange}{HTML}{FF6D00}

\setbeamercolor{structure}{fg=MyDarkTeal}

\setbeamercolor{headline}{bg=MyDarkTeal, fg=white}
\setbeamercolor{footline}{bg=MyDarkTeal, fg=white}

\setbeamercolor{frametitle}{bg=MyDarkTeal, fg=white}

\setbeamercolor{block title}{bg=MyDarkTeal, fg=white}
\setbeamercolor{block body}{bg=MySoftPink!25, fg=black}

\setbeamercolor{block title alerted}{bg=MyVividOrange, fg=white}
\setbeamercolor{block body alerted}{bg=MyVividOrange!15, fg=black}

\setbeamercolor{block title example}{bg=MySoftPink, fg=MyDarkTeal}
\setbeamercolor{block body example}{bg=MySoftPink!15, fg=black}

\setbeamercolor{alerted text}{fg=MyVividOrange}
\setbeamercolor{item}{fg=MyVividOrange}

% Operators
\DeclareMathOperator{\Spec}{Spec}
\DeclareMathOperator{\rad}{rad}
\DeclareMathOperator{\Div}{Div}
\DeclareMathOperator{\CaDiv}{CaDiv}
\DeclareMathOperator{\Cl}{Cl}
\DeclareMathOperator{\CaCl}{CaCl}
\DeclareMathOperator{\Pic}{Pic}
\DeclareMathOperator{\divv}{div}
\DeclareMathOperator{\Hom}{Hom}
\DeclareMathOperator{\im}{im}
\DeclareMathOperator{\coker}{coker}
\newcommand{\OO}{\mathcal{O}}
\newcommand{\FF}{\mathcal{F}}
\newcommand{\GG}{\mathcal{G}}
\newcommand{\KK}{\mathcal{K}}
\newcommand{\LL}{\mathcal{L}}
\newcommand{\A}{\mathbb{A}}
\newcommand{\PP}{\mathbb{P}}
\newcommand{\Z}{\mathbb{Z}}
\newcommand{\pp}{\mathfrak{p}}
\newcommand{\mm}{\mathfrak{m}}
\newcommand{\aaa}{\mathfrak{a}}

\theoremstyle{definition}
\newtheorem{proposition}[theorem]{Proposition}
\newtheorem{remark}[theorem]{Remark}

\title[Divisor of Algebraic Curves\texorpdfstring{\quad\insertframenumber/\inserttotalframenumber}{}]{Divisor of Algebraic Curves (Foundations Ver.)}
\subtitle{
Tugas Seminar I version \\    
MA4091 Seminar Matematika I}
\author[G. Allister (10123021)]{Gilbert Allister (10123021) \\
Dosen Pembimbing: Rizal Afgani, Ph.D. \\
Dosen Penguji: 
}


\date{ \fontsize{8pt}{8pt}\selectfont
  Program Studi Sarjana Matematika \\
  Fakultas Matematika dan Ilmu Pengetahuan Alam \\
  Institut Teknologi Bandung \\[1cm]
  \today} % presentation date (was \today)

\begin{document}
    \frame{\titlepage}

    \begin{frame}{Table of Contents}
        \tableofcontents
    \end{frame}

    \section{Introduction}

    \begin{frame}{Description of this Presentation}
        \begin{itemize}
            \item The Topic of this Final Paper is "Divisors of Algebraic Curves".
            \item The foundations of this topic are Classical Algebraic Geometry, Scheme Theory and Commutative Algebra.
            \item While the key machinaries of the topic are Sheaf Theory, Invertible Sheaves, and Weil/Cartier Divisors. 
            \item Hence in this presentation, we shall present some foundational topics.
        \end{itemize}
    \end{frame}


    \begin{frame}{Motivation}
        \begin{itemize}
            \item Given polynomials $p_1, \dots, p_r \in k[X_1, \dots, X_n]$, we wish to understand the geometry of the solutions of $p_1 = \dots = p_r = 0$.
            \item For example, $p(x,y) = x^2 + y^2 - 1$ cuts out a circle: one equation in the plane gives a \emph{curve}.
            \item On a curve we can ask for the zeros and poles of a rational function, \emph{together with their multiplicities}.
            \item For example, $f = \dfrac{x-1}{x^2}$ on $\A^1_k$ has a simple zero at $1$ and a double pole at $0$. We record this as the formal sum $[1] - 2[0]$.
            \item Such formal sums of points are called \alert{divisors}. To make this precise, we need rings, topology and sheaves.
        \end{itemize}
    \end{frame}

    \begin{frame}{Roadmap}
        \begin{enumerate}
            \item \textbf{Noetherian rings, Zariski topology, Nullstellensatz:} the dictionary between algebra and geometry.
            \item \textbf{Affine varieties:} coordinate rings, polynomial maps, rational functions, local rings.
            \item \textbf{Sheaves and affine schemes:} gluing local data; $\Spec A$ and its structure sheaf.
            \item \textbf{Looking ahead:} Weil divisors, Cartier divisors, invertible sheaves and $\Pic$ (the main body of the paper).
        \end{enumerate}
    \end{frame}

    \section{Noetherian, Zariski and Nullstellesatz}

    \begin{frame}{Outline of this section}
        \begin{enumerate}
            \item Noetherian Rings
            \item Zariski Topology
            \item Nullstellensatz Theorem
        \end{enumerate}
        
    \end{frame}

    \begin{frame}{Noetherian Rings}

        \begin{definition}[Noetherian Ring]
            A ring $A$ is \emph{Noetherian} if the following equivalent conditions hold:
            \begin{enumerate}
                \item[(i)] Every ideal $I \subset A$ is finitely generated; that is, for every ideal $I \subset A$, there exist $f_1, \dots, f_k \in I$ such that $I = (f_1, \dots, f_k)$.
                \item[(ii)] Every ascending chain
                \begin{align*}
                    I_1 \subset I_2 \subset \dots \subset I_m \subset \dots
                \end{align*}
                of ideals of $A$ terminates, that is, the chain is eventually stationary, with $I_N = I_{N+1} = \dots$ (the ascending chain condition, or a.c.c.).
                \item[(iii)] Every nonempty set of ideals of $A$ has a maximal element.
            \end{enumerate}
        \end{definition}

        \begin{itemize}
            \item For example, $\Z$ and $k[x]$ are Noetherian.
            \item The intuition of Noetherian Ring is we want to work in ring that has finite generators.
        \end{itemize}

    \end{frame}


    \begin{frame}
        \begin{proposition}
            (i) Suppose that $A$ is Noetherian, and $I \subset A$ an ideal; then the quotient ring $B = A/I$ is Noetherian.

            (ii) Let $A$ be a Noetherian integral domain, and $A \subset K$ its field of fractions; let $\emptyset \neq S \subset A$ be a subset, and set
            \begin{align*}
                B = A[S^{-1}] = \left\{\frac{a}{b} \in K \;\middle|\; a \in A,\; b = 1\text{ or a product of elements of }S\right\}.
            \end{align*}
            Then $B$ is again Noetherian.
        \end{proposition}

        \begin{theorem}[Hilbert Basis Theorem] \label{thm:hilbert-basis}
            For a ring $A$,
            \begin{align*}
                A \text{ Noetherian} \Rightarrow A[X] \text{ Noetherian}.
            \end{align*}
        \end{theorem}

        \begin{corollary}
            For $k$ a field, a finitely generated $k$-algebra is Noetherian.
        \end{corollary}
    \end{frame}

    \begin{frame}{Algebraic set}
        \begin{definition} \label{def:algebraic-set}
            A subset $X \subset \mathbb{A}^n_k$ is an \emph{algebraic set} if $X = V(I)$ for some ideal $I$. Notice that by Corollary 3.3, $I$ is finitely generated. If $I = (f_1, \dots, f_r)$ then clearly
            \begin{align*}
                V(I) = \{P \in \mathbb{A}^n_k \mid f_i(P) = 0 \text{ for } i = 1, \dots, r\},
            \end{align*}
        so that an algebraic set is just a locus of points satisfying a finite number of polynomial equations.
        \end{definition}
    \end{frame}

    \begin{frame}{Zariski Topology}
        \begin{proposition}[Proposition-Definition] \label{prop:zariski}
            The correspondence $V$ satisfies the following formal properties:
            \begin{enumerate}
            \item[(i)] $V(0) = \mathbb{A}^n_k$;\quad $V(A) = \emptyset$;
            \item[(ii)] $I \subset J \implies V(J) \subset V(I)$;
            \item[(iii)] $V(I_1 \cap I_2) = V(I_1) \cup V(I_2)$;
            \item[(iv)] $V\left(\sum_i I_i\right) = \bigcap_i V(I_i)$.
            \end{enumerate}
            Hence the algebraic subsets of $\mathbb{A}^n_k$ form the closed sets of a topology on $\mathbb{A}^n_k$, the \emph{Zariski topology}.
        \end{proposition}
    \end{frame}


    \begin{frame}{V-I Correspondence}

        Let $A = k[X_1, \dots, X_n]$.
        \begin{align*}
            \{\text{ideals } J \subset A\} \longrightarrow \{\text{subsets } X \subset \mathbb{A}^n_k\}  \\
            J \mapsto V(J) = \{P \in \mathbb{A}^n_k \mid f(P) = 0 \text{ for all } f \in J\}.
        \end{align*}

        As a kind of inverse to $V$ there is a correspondence
        \begin{align*}
            \{\text{ideals } J \subset A\} &\longleftarrow \{\text{subsets } X \subset \mathbb{A}^n_k\} \\
            I(X) = \{f \in k[X_1, \dots, X_n] \mid f(P) = 0 \text{ for all } P \in X\}. &\mapsfrom X
        \end{align*}
        That is, $I$ takes a subset $X$ to the ideal of functions vanishing on it.
    \end{frame}

    \begin{frame}{Radical Ideal}
        \begin{itemize}
            \item Before stating the Nullstellensatz, we define the radical of an ideal in order to simplify the statement.
            \begin{definition}[Radical ideal]
                Let $I$ be an ideal of $A = k[X_1, \ldots, X_n]$. The \textbf{radical} of $I$ is
                \begin{align*}
                    \rad I = \{ f \in A : f^m \in I \text{ for some integer } m \geq 1\}.
                \end{align*}
                If $I = \rad I$, we say that $I$ is a \textbf{radical ideal}.
            \end{definition}
            \item $\rad I$ is an ideal containing $I$, and $V(\rad I) = V(I)$.
            \item Every prime ideal is radical, and $I(X)$ is a radical ideal for every $X \subseteq k^n$.
            \item For example, in $k[X]$ we have $\rad\,(X^2) = (X)$. Both ideals define the same algebraic set $\{0\}$, so the algebraic set does not remember the multiplicity.
        \end{itemize}
    \end{frame}

    \begin{frame}{Hilbert's Nullstellensatz}
        \begin{theorem}[Nullstellensatz, weak form]
            Let $k$ be an algebraically closed field and $A = k[X_1, \ldots, X_n]$. If $\mm$ is a maximal ideal of $A$, then
            \begin{align*}
                \mm = (X_1 - a_1,\, X_2 - a_2,\, \ldots,\, X_n-a_n)
            \end{align*}
            for some point $(a_1, a_2, \ldots, a_n) \in k^n$.
        \end{theorem}

        \begin{theorem}[Nullstellensatz, strong form]
            Let $k$ be an algebraically closed field and $A = k[X_1, \ldots, X_n]$. If $J$ is an ideal of $A$ with $J \neq A$, then $V(J) \neq \emptyset$. Furthermore, for every ideal $J$ of $A$,
            \begin{align*}
                I(V(J)) = \rad J.
            \end{align*}
        \end{theorem}
    \end{frame}

    \begin{frame}{The Algebra--Geometry Dictionary}
        \begin{corollary}
            The maps $X \mapsto I(X)$ and $J \mapsto V(J)$ are mutually inverse, inclusion-reversing bijections:
            \begin{center}
                \upshape
                \begin{tabular}{rcl}
                    \toprule
                    \textbf{Geometry} (in $k^n$) & & \textbf{Algebra} (in $A = k[X_1,\ldots,X_n]$) \\
                    \midrule
                    algebraic sets & $\longleftrightarrow$ & radical ideals \\
                    irreducible & $\longleftrightarrow$ & prime ideals \\
                    points & $\longleftrightarrow$ & maximal ideals \\
                    \bottomrule
                \end{tabular}
            \end{center}
        \end{corollary}
        \begin{itemize}
            \item So the geometry of $k^n$ can be recovered from the ring $A$ alone.
            \item \textbf{Idea of scheme theory:} replace $A$ by an arbitrary commutative ring with $1$, and take \emph{all} prime ideals as the points of the space.
        \end{itemize}
    \end{frame}

    \section{Affine Varieties}

    \begin{frame}{Outline of this section}
        
    \end{frame}

    \begin{frame}{Polynomial Maps}
        \begin{definition}{Polynomial Maps}
            Define $f \colon V \to W$ as polynomial maps, given by $m$ polynomial functions $f_1, \dots, f_m \in k[V],$, so that for $P \in V$,
            \begin{align*}
                f(P) = \bigl(f_1(P), \dots, f_m(P)\bigr) \in W.
            \end{align*}
        \end{definition}

        \begin{itemize}
            \item Keep in mind that if we have composite of $f$  with any polynomial function $g \in k[W]$ is again polynomial function on $V$. 
            \item We can see that polynomial maps evaluate each coordinates by its polynomial functions, hence we have such proposition.
        \end{itemize}

        \begin{proposition}
            Suppose we have $V \subset \A^n_k$, $W \subset \A^m_k$, and we write  $f = (f_1, \ldots, f_m)$, then 
            \begin{align*}
                f: V \to W \text{ is polynomial map} \iff \forall j, f_j = Y_j \circ f \in k[W] 
            \end{align*}
            by given $Y_j$ $j$-th coordinate function.
        \end{proposition}
        
    \end{frame}


    \begin{frame}
        \begin{itemize}
            \item Given polynomial map $f: V \to W$, now we define map 
                \begin{align*}
                    f^* : k[W] &\to k[V] \\
                    g &\mapsto g \circ f
                \end{align*}
            \item Then $f^*$ is ring homomorphism, and in fact, it is k-algebra homomorphism.
            \item Note that it went backwards: from coordinate ring of the target to that of the source.
            \item Now with that, we propose a claim 
            \begin{proposition}
                Every $k$-algebra homomorphism $\Phi \colon k[W] \to k[V]$ is of the form $\Phi = f^*$ for a \emph{uniquely defined} polynomial map $f \colon V \to W$.
            \end{proposition}
        \end{itemize}

        
        
    \end{frame}

    \begin{frame}
        \begin{itemize}
            \item Before defining affine variety, we want the following claim.
            
            \begin{proposition}
                Given $V \subset \A^n_k$, we have 
                \begin{align*}
                    V \text{ is irreducible} \iff I(V) \text{ prime ideal} \iff k[V] \text{ is integral domain}
                \end{align*}
            \end{proposition}
            \textit{Sketch proof.} We shall prove the contrapositive $X$ reducible $\leftrightarrow$ $I(X)$ is not prime. If $X$ is  reducible, then we have $X = X_1 \cup X_2$ and choose $f_1 = I(X_1)\backslash I(X)$ and $f_2 = I(X_2) \backslash I(X)$, hence $f_1, f_2 \notin I(X)$ but $f_1f_2 \in I(X)$ since it vanishes there. If the ideal is not prime, then $\exists f_1, f_2 \notin I(X)$ but $f_1f_2 \in I(X)$. Construct $I_1 = (I(X), f_1)$ and $X_1 = V(I_1)$, while samewise with $X_2 = V(I_2)$, then $X = X_1 \cup X_2$. The second iff is obvious since $k[V] = k[X_1, \ldots, X_n]/I(V)$.      
            
            \item We want $V$ be irreducible, so that we can use $k[V]$ to define a quotient field, which will be shown later.
        \end{itemize}
        
    \end{frame}

    \begin{frame}{Affine Varieties}

        \begin{definition}[Affine varieties]
            An \emph{affine variety} over a field $k$ is an irreducible algebraic subset $V \subset \mathbb{A}^n_k$, equipped with its coordinate ring $k[V]$. Equivalently, it is a set $V$ together with a ring $k[V]$ of $k$-valued functions $f\colon V \to k$ such that:
            \begin{enumerate}
                \item $k[V]$ is a finitely generated $k$-algebra, and
                \item for some choice $x_1, \dots, x_n$ of generators of $k[V]$ over $k$, the map
                \begin{align*}
                    V \to \mathbb{A}^n_k, \qquad P \mapsto \bigl(x_1(P), \dots, x_n(P)\bigr)
                \end{align*}
                embeds $V$ as an irreducible algebraic set.
            \end{enumerate}
        \end{definition}
    \end{frame}

    \begin{frame}{Morphism between Affine Varieties}
        \begin{itemize}
            \item Now we wish to know the mapping between affine varieties. 
            \item Intuitively, we would use usual mapping consists of polynomial function, as previous shown.
            \item However, it is too restrictive as we can have rational function.
            \item For example, suppose we have Pythogoras equation $X^2 + Y^2 = Z^2$ and we wish to transform it into circle by $x = X/Z$ and $y = Y/Z$, so we have $x^2 + y^2 = 1$ so finally we can have parametrisation of $\lambda = \frac{x}{1-y}$, hence $x = \frac{2 \lambda}{\lambda^2 + 1}$ and $y = \frac{\lambda^2 - 1}{\lambda^2 + 1}$   
            \item And notice that using the mapping of $f(x) = 1/x$ is not polynomial, hence we need to define rational mapping.
        \end{itemize}
        
    \end{frame}

    \begin{frame}{Rational Function}
        \begin{itemize}
            \item Given $V$ is irreducible and $k[V]$ domain integral, then define 
            \begin{align*}
                k(V) := \text{Quot}(k[V])
            \end{align*}  
            as $\text{Quot}(R)$ is quotient field of domain integral $R$. 

            \item By this, we can write $f \in k(V)$ as $f = g/h$ for $h \neq 0$ (in the sense of zero polynomial).  
            \item But notice $f$ is not well defined if $\exists P  \in V$ such that $h(P) = 0$.    
            \item Hence we define as following 
            \begin{definition}
                Let $f \in k(V)$ and $P \in V$, define $f$ is \emph{regular} at $P$ if $h(P) \neq 0$ with $f = g/h$, $g,h \in k[V]$, and define \emph{domain of definition} of $f$ by 
                \begin{align*}
                    \operatorname{dom} f := \{ P \in V \hspace{3pt} | \hspace{3pt} f \text{ is regular at } P \}
                \end{align*}   
                and define \emph{local ring} of $V$ at $P$ as 
                \begin{align*}
                    \OO_{V, P} := \{ f \in k(V) \hspace{3pt} | \hspace{3pt} f \text{ is regular at } P \}
                \end{align*}      
            \end{definition}
            
        \end{itemize}

        
    \end{frame}

    \begin{frame}
        \begin{itemize}
            \item We seek the characteristics of the $\operatorname{dom} f$, define $D_f := \{ h \in k[V] \hspace{3pt} | \hspace{3pt} hf \in k[V] \} $, then
            \begin{align*}
                V \backslash \operatorname{dom} f = \{ P \in V \hspace{3pt} | \hspace{3pt} h(P) = 0 \} = V(D_f)
            \end{align*}
            then $\operatorname{dom} f $ open in Zariski Topology. Since it is not empty, then it is dense.
            \item $\operatorname{dom} f = V$ if and only if $f \in k[V]$, i.e.\ $f$ is a polynomial function, since 
            \begin{align*}
                \operatorname{dom} f = V \iff V(D_f) = \emptyset \iff 1 \in D_f \iff f \in k[V] 
            \end{align*}
            \item For any $h \in k[V]$, let $V_h = V \setminus V(h) = \{ P \in V \mid h(P) \neq 0\}$. Then
            \begin{align*}
                \operatorname{dom} f \supset V_h \quad\Longleftrightarrow\quad f \in k[V][h^{-1}].
            \end{align*}
            since 
            \begin{align*}
                \operatorname{dom} f \supset V_h \iff h \text{ vanishes on } D_f \iff h^n \in D_f \text{ for some n} \iff f = g/h^n \in k[V][h^{-1}]
            \end{align*}
        \end{itemize}
        
    \end{frame}

    \begin{frame}
        \begin{itemize}
            
            \item Hence, we can write the local ring as such 
            \begin{align*}
                \OO_{V, P} = k[V][(h^{-1} \hspace{3pt} | \hspace{3pt} h(P) \neq 0)]
            \end{align*}
            \item Finally, we can define rational map 
            \begin{definition}
                A \emph{rational map} $f \colon V \dashrightarrow \mathbb{A}^m_k$ is a partially defined map given by rational functions $f_1, \dots, f_n \in k(V)$, that is,
                \begin{align*}
                    f(P) = \bigl(f_1(P), \dots, f_m(P)\bigr) \quad \text{for all } P \in \bigcap_i \operatorname{dom} f_i.
                \end{align*}
                By definition, $\operatorname{dom} f = \bigcap_i \operatorname{dom} f_i$. $f$ is said to be \emph{regular} at $P \in V$ if and only if $P \in \operatorname{dom} f$.

                A rational map $f \colon V \dashrightarrow W$ between two affine varieties $V \subset \mathbb{A}^n_k$ and $W \subset \mathbb{A}^m_k$ is defined to be a rational map $f \colon V \dashrightarrow \mathbb{A}^m_k$ such that $f(\operatorname{dom} f) \subset W$.
            \end{definition}
        \end{itemize}
        
    \end{frame}

    \begin{frame}
        \begin{itemize}
            \item For example, consider $C = V(x^2 + y^2 - 1)$. The maps
            \begin{align*}
                \A^1 \dashrightarrow C,\ \lambda \mapsto \left(\tfrac{2\lambda}{\lambda^2+1}, \tfrac{\lambda^2-1}{\lambda^2+1}\right), \qquad
                C \dashrightarrow \A^1,\ (x,y) \mapsto \tfrac{x}{1-y}
            \end{align*}
            although the maps are not polynomial (the first is undefined where $\lambda^2 + 1 = 0$, the second where $y = 1$).
        \end{itemize}
        
    \end{frame}



    \section{Sheaf \text{\&} Affine Scheme}

    \begin{frame}{Outline of this section}
        
    \end{frame}

    \begin{frame}{Why Sheaves?}
        \begin{itemize}
            \item Now we back to Nullstellensatz, where the idea is to replace the $A$ by an arbitrary commutative ring with $1$, and take \emph{all} prime ideals as the points of the space. 
            \item Before that, we need to define several objects.
            
            \begin{definition}[Presheaf]
            Let $X$ be a topological space. A \emph{presheaf} $\FF$ of rings on $X$ assigns
            \begin{itemize}
                \item to each open $U \subset X$ a ring $\FF(U)$ (the \emph{sections} over $U$), and
                \item to each inclusion $V \subset U$ a ring homomorphism $\rho_{UV} \colon \FF(U) \to \FF(V)$, written $s \mapsto s|_V$,
            \end{itemize}
            such that $\FF(\emptyset) = 0$, $\rho_{UU} = \text{id}$, and $\rho_{VW} \circ \rho_{UV} = \rho_{UW}$ for $W \subset V \subset U$.
        \end{definition}
        \end{itemize}
        
    \end{frame}

    \begin{frame}{Sheaf Axioms}
        \begin{itemize}
            \item Then we shall begin to define sheaf 
            \begin{definition}[Sheaf]
            A presheaf $\FF$ on $X$ is a \emph{sheaf} if for every open $U$ and every open cover $U = \bigcup_i U_i$:
            \begin{enumerate}
                \item[(S1)] \emph{(Locality)} if $s \in \FF(U)$ and $s|_{U_i} = 0$ for all $i$, then $s = 0$;
                \item[(S2)] \emph{(Gluing)} if $s_i \in \FF(U_i)$ satisfy $s_i|_{U_i \cap U_j} = s_j|_{U_i \cap U_j}$ for all $i, j$, then there is $s \in \FF(U)$ with $s|_{U_i} = s_i$ for all $i$.
            \end{enumerate}
            (By S1, the section $s$ in S2 is unique.)
        \end{definition}
            \item For examples, Continuous real functions on a topological space $X$ form a sheaf.
            \item And related example, for an affine variety $V$, $U \mapsto \OO_V(U)$ is a sheaf, the \emph{sheaf of regular functions}.
            \item For non-example, the \emph{constant} presheaf $U \mapsto \mathbb{R}$ is not a sheaf: if $U = U_1 \sqcup U_2$ is disconnected, the sections $0$ on $U_1$ and $1$ on $U_2$ agree on the (empty) overlap but do not glue to a constant on $U$.
        \end{itemize}
    \end{frame}

    \begin{frame}{Stalks}
        \begin{definition}[Stalk]
            For a presheaf $\FF$ on $X$ and $x \in X$, the \emph{stalk} at $x$ is
            \begin{align*}
                \FF_x := \varinjlim_{U \ni x} \FF(U),
            \end{align*}
            that is, germs of sections: $s \in \FF(U)$ and $t \in \FF(U')$ define the same germ iff $s|_W = t|_W$ for some open $W \ni x$ inside $U \cap U'$.
        \end{definition}
        \begin{itemize}
            \item For the sheaf $\OO_V$ of an affine variety, the stalk at $P \in V$ is exactly the local ring $\OO_{V,P}$ defined earlier.
            \item Stalks are where we do ``local'' algebra; they are what will let us define $\mathrm{ord}_P$ on a curve.
        \end{itemize}
    \end{frame}

    \begin{frame}{The Spectrum of a Ring}
        \begin{definition}[$\Spec A$]
            Let $A$ be a commutative ring with $1$. As a set,
            \begin{align*}
                \Spec A := \{ \pp \subset A \mid \pp \text{ is a prime ideal} \}.
            \end{align*}
            For an ideal $I \subset A$ put $V(I) = \{ \pp \in \Spec A \mid \pp \supseteq I \}$.
        \end{definition}
        \begin{itemize}
            \item $V(I)$ satisfies the same four properties as in the Zariski Topology slide, so the $V(I)$ are the closed sets of a topology on $\Spec A$ (the \emph{Zariski topology}).
            \item For $f \in A$, the \emph{principal open sets}
            \begin{align*}
                D(f) := \Spec A \setminus V(f) = \{ \pp \mid f \notin \pp \}
            \end{align*}
            form a basis of the topology.
            \item $\{\pp\}$ is closed iff $\pp$ is maximal; in general $\overline{\{\pp\}} = V(\pp)$. If $A$ is a domain, $(0)$ is a \emph{generic point}: $\overline{\{(0)\}} = \Spec A$.
        \end{itemize}
    \end{frame}

    \section{Closing \text{\&} References}

    \begin{frame}{References}
        \begin{enumerate}
            \item M. Reid, \emph{Undergraduate Algebraic Geometry}, London Mathematical Society Student Texts 12, Cambridge University Press, 1988.
            \item U. G\"ortz and T. Wedhorn, \emph{Algebraic Geometry I: Schemes with Examples and Exercises}, Springer Vieweg, 2010.
            \item 
        \end{enumerate}
        
    \end{frame}

\end{document}
